\chapter{Equiangular Frames and Grassmannian Packings}
\label{ch:etf}
\fancyhead[R]{\small\textit{Chapter 9: Equiangular and Grassmannian Frames}}
\section{Overview: Beyond Tight, Toward Optimal}
\label{sec:ch9-overview}

Chapter~8 showed that tight, unit-norm frames are the global minimizers
of the frame potential --- the most ``spread out'' unit-norm
configurations in the energy-minimization sense. Lab Exercise~8.4 asked
a natural follow-up question: among those minimizers, are the resulting
vectors \emph{equiangular} (all pairwise inner product magnitudes equal),
or merely tight? The answer, observed numerically for $m=7$, $n=3$, was:
merely tight --- gradient descent on the frame potential finds tight
frames, but not necessarily equiangular ones, because equiangularity is
a strictly stronger condition than tightness.

This chapter studies that stronger condition. An \emph{equiangular tight
frame} (ETF) is a unit-norm tight frame in which \emph{every pair} of
distinct frame vectors has exactly the same inner product magnitude.
ETFs are the most ``spread out'' configurations in a second, independent
sense: they minimize the maximum pairwise inner product over all
unit-norm frames of the same size and ambient dimension. The precise
lower bound on that maximum is the \emph{Welch bound}, and ETFs are
exactly the configurations that achieve it.

We also revisit the connection, promised in Chapter~4, between the
roots-of-unity construction and a larger family called \emph{harmonic
frames}. The chapter closes with the \emph{Grassmannian packing
problem}: the question of how to arrange $m$ one-dimensional subspaces
in $\R^n$ (or $\C^n$) as spread out as possible, of which ETFs are
one specific and very clean solution.

\section{Coherence and Equiangular Frames}
\label{sec:coherence}

\begin{definition}[Coherence]
\label{def:coherence}
Let $\{f_i\}_{i=1}^m \subset \F^n$ be a collection of unit-norm
vectors. The \emph{coherence} of the collection is
\[
  \mu\bigl(\{f_i\}\bigr) \;:=\;
  \max_{1\le i < j \le m} |\inner{f_i}{f_j}|.
\]
\end{definition}

The coherence is the worst-case pairwise inner product magnitude: it
measures how ``similar'' the two most similar vectors in the collection
are. Small coherence means no two vectors are nearly parallel --- the
configuration is well spread out in the sense of
pairwise-distinguishability. A set of unit-norm vectors with
$\mu=0$ is an orthonormal set (every pair is orthogonal); increasing
$\mu$ indicates growing redundancy, with $\mu=1$ achieved when some pair
of vectors is actually parallel (or antipodal, since $|\inner{f_i}{-f_i}|=1$).

\begin{remark*}
Coherence is not the same as the frame's lower bound $A$ or upper bound
$B$: coherence is a property of \emph{pairs} of vectors, while frame
bounds are global properties of the entire configuration. We have been
using frame bounds since Chapter~2; coherence is a finer, more local
notion that appears naturally when asking how well a frame can
distinguish nearly-parallel directions.
\end{remark*}

\begin{definition}[Equiangular frame]
\label{def:equiangular}
A collection $\{f_i\}_{i=1}^m$ of unit-norm vectors in $\F^n$ is
\emph{equiangular} if there is a constant $c\ge0$ such that
$|\inner{f_i}{f_j}|=c$ for every $i\ne j$. Equivalently, all off-diagonal
entries of the Gram matrix $G_{ij}=\inner{f_j}{f_i}$ have the same
modulus. When the collection is also a tight frame, it is called an
\emph{equiangular tight frame} (ETF).
\end{definition}

\begin{examplebox}[title={Example 9.1: The Mercedes-Benz frame}]
\label{ex:mercedes-benz}
The three vectors $f_1=(1,0)$, $f_2=(-\tfrac12,\tfrac{\sqrt3}{2})$,
$f_3=(-\tfrac12,-\tfrac{\sqrt3}{2})$ in $\R^2$ --- the triangular frame
studied since Chapter~1 --- are the canonical example of an ETF. They
satisfy $|\inner{f_i}{f_j}|=\tfrac12$ for every $i\ne j$ (all
off-diagonal inner products equal $-\tfrac12$, which in the real case
gives modulus $\tfrac12$), and are tight with $A=\tfrac32$
(Example~3.2). This frame is sometimes called the \emph{Mercedes-Benz
frame} after the visual resemblance of its three vectors to the three
spokes of the well-known logo.
\end{examplebox}

\begin{remark*}
Not every tight frame is equiangular. The square frame ($m=4$, $n=2$,
the vertices of a square) is tight with $A=2$ (Chapter~4, Example~3.6)
but is manifestly \emph{not} equiangular: adjacent vertices have inner
product $0$, while antipodal vertices have inner product $|\inner{f_1}{f_3}|
= |(1,0)\cdot(-1,0)| = 1$. The off-diagonal inner products take two
distinct values, $0$ and $1$, so $\mu = 1$ --- the worst possible
coherence for a unit-norm set. Any regular $m$-gon with $m\ge4$ has
this deficiency: if $m$ is even, antipodal pairs give coherence $1$;
if $m$ is odd with $m\ge5$, at least two distinct pairwise angles appear
(as verified numerically in Lab Exercise~9.1). Among regular polygons
in $\R^2$, \emph{only the triangle} ($m=3$) is equiangular, making the
Mercedes-Benz frame the unique real ETF arising from the regular-polygon
family.
\end{remark*}

\section{The Welch Bound}
\label{sec:welch-bound}

Why does it matter whether the coherence equals some specific value?
The \emph{Welch bound} provides a precise, universal lower limit on the
coherence of any unit-norm collection, expressed in terms of $m$ and $n$
alone. It says you cannot spread $m$ unit-norm vectors in $\F^n$ with
coherence below a certain threshold --- no matter how hard you try.

\begin{theorem}[Welch bound]
\label{thm:welch-bound}
For any collection $\{f_i\}_{i=1}^m$ of unit-norm vectors in $\F^n$
with $m>n$,
\[
  \mu\bigl(\{f_i\}\bigr)
  \;\ge\;
  \sqrt{\frac{m-n}{n(m-1)}},
\]
with equality if and only if $\{f_i\}_{i=1}^m$ is an equiangular tight
frame.
\end{theorem}

\begin{proof}
Let $G$ be the $m\times m$ Gram matrix with entries $G_{ij}=
\inner{f_j}{f_i}$. Since all $f_i$ are unit-norm, the diagonal entries
are $G_{ii}=1$, so $\tr(G)=m$. The off-diagonal entries satisfy
$|G_{ij}|\le\mu$ for $i\ne j$.

\textbf{Step 1: lower-bound $\tr(G^2)$.}
Since $G=F^*F$ and $S=FF^*$ are related by $G=(F^*)(F)$ and $S=(F)(F^*)$,
they share the same nonzero eigenvalues (a standard fact: if $ABv=\lambda v$
then $BA(Bv)=\lambda(Bv)$, so nonzero eigenvalues of $AB$ and $BA$ coincide).
Therefore $\tr(G^2)=\sum_{k}\sigma_k^2$ where $\sigma_k$ are the eigenvalues
of $G$, which equal those of $S$ for the $n$ nonzero ones (and $0$ for the
remaining $m-n$). In particular, $\tr(G^2)=\tr(S^2)=\FP(\{f_i\})$. By
Theorem~8.3.1 (frame potential lower bound), $\FP\ge m^2/n$ for any
unit-norm collection, giving
\[
  \tr(G^2) \;\ge\; \frac{m^2}{n}.
\]

\textbf{Step 2: upper-bound $\tr(G^2)$ using $\mu$.} Split the sum:
\[
  \tr(G^2)=\sum_{i,j}|G_{ij}|^2
  = \underbrace{\sum_i |G_{ii}|^2}_{=m}
    + \underbrace{\sum_{i\ne j}|G_{ij}|^2}_{\le\,m(m-1)\mu^2}.
\]
So $\tr(G^2)\le m + m(m-1)\mu^2$.

\textbf{Step 3: combine.} From Steps~1 and~2:
\[
  \frac{m^2}{n} \;\le\; \tr(G^2) \;\le\; m + m(m-1)\mu^2.
\]
Rearranging:
\[
  m(m-1)\mu^2 \;\ge\; \frac{m^2}{n} - m = m\!\left(\frac{m}{n}-1\right)
  = \frac{m(m-n)}{n},
\]
so $\mu^2 \ge (m-n)/(n(m-1))$, which gives the Welch bound.

\textbf{Equality.} The two Cauchy--Schwarz applications give equality
simultaneously iff: (a) all eigenvalues of $S$ are equal, i.e.\ $S=\lambda
I$ (tight frame, by Theorem~3.5.1(iii)); and (b) all off-diagonal
$|G_{ij}|$ are equal to $\mu$ (equiangular). Together, these say
$\{f_i\}$ is an ETF.
\end{proof}

\begin{keyidea}
\textbf{ETFs are exactly the configurations achieving the Welch bound}
--- the configurations of unit-norm vectors whose worst-case pairwise
inner product is as small as it can possibly be, given $m$ and $n$. The
Welch bound is a hard floor: no unit-norm configuration of $m$ vectors
in $\F^n$ can have coherence below $\sqrt{(m-n)/(n(m-1))}$. ETFs sit on
that floor. Tight frames that are not ETFs (e.g.\ regular polygons with
$m\ge4$ in $\R^2$) sit strictly above it.
\end{keyidea}

\begin{examplebox}[title={Example 9.2: Welch bound for the Mercedes-Benz frame}]
\label{ex:welch-mercedes}
For $m=3$, $n=2$: Welch bound $= \sqrt{(3-2)/(2\cdot2)} = \sqrt{1/4}
= \tfrac12$. The Mercedes-Benz frame achieves $\mu=\tfrac12$ exactly
(Example~\ref{ex:mercedes-benz}). So it sits precisely on the Welch
bound --- it is the tightest possible unit-norm configuration of $3$
vectors in $\R^2$ in the coherence sense, among all possible
configurations of $3$ unit-norm vectors in $\R^2$.

For comparison, the square frame ($m=4$, $n=2$) has Welch bound
$\sqrt{(4-2)/(2\cdot3)} = \sqrt{1/3} \approx 0.577$, but achieves
$\mu=1$ (antipodal pair). It sits far above the Welch floor, confirming
that being tight does not imply achieving the Welch bound.
\end{examplebox}

\begin{remark*}
The Welch bound is achievable only for certain pairs $(n,m)$ --- there
are combinatorial and algebraic obstructions that prevent ETFs from
existing in many cases. For example, in $\R^2$, the only real ETF with
$m>n$ is the Mercedes-Benz frame ($m=3$): no unit-norm equiangular tight
frame with $m\ge4$ vectors exists in $\R^2$. (This can be verified
directly: equiangularity forces all $|\inner{f_i}{f_j}|=c$ for a common
$c>0$, tightness forces $c=\sqrt{m-n}/(n(m-1))^{1/2}$ which is the
Welch bound, and in $\R^2$ the only unit-norm vectors satisfying both
constraints simultaneously for $m\ge4$ would have to include antipodal
pairs, which gives $c=1$, contradicting $c<1$ for $m>n+1$... the precise
combinatorial statement is subtle, but the numerical evidence is clear.)
In higher dimensions and over $\C$, ETFs can exist for many more
$(n,m)$ pairs, and their construction is an active research area
connecting to combinatorics, coding theory, and quantum information.
\end{remark*}

\section{An ETF Characterization via the Gram Matrix}
\label{sec:gram-characterization}

The definition of an ETF involves two separate conditions: tight and
equiangular. There is an elegant way to unify them into a single
condition on the Gram matrix, which also gives a concrete necessary
condition for ETFs to exist.

\begin{proposition}[ETF Gram matrix characterization]
\label{prop:etf-gram}
A unit-norm collection $\{f_i\}_{i=1}^m \subset \F^n$ is an ETF with
bound $A=m/n$ and coherence $c$ if and only if its Gram matrix
$G$ satisfies
\begin{equation}
\label{eq:etf-gram}
  G^2 = \frac{m}{n}\, G.
\end{equation}
\end{proposition}

\begin{proof}
Since $G=F^*F$ (where $F$ has columns $f_i$) and $S=FF^*$, we have
$G^2=(F^*F)^2=F^*(FF^*)F=F^*SF$. The frame is tight with bound $A=m/n$
iff $S=(m/n)I$, in which case $F^*SF=(m/n)F^*F=(m/n)G$, giving
$G^2=(m/n)G$. Conversely, if $G^2=(m/n)G$, then $G$ is a projection
(up to scale $m/n$), meaning the eigenvalues of $G$ lie in $\{0, m/n\}$.
Since $\tr(G)=m$ and $G$ has $m$ eigenvalues summing to $m$, exactly $n$
eigenvalues equal $m/n$ and the remaining $m-n$ equal $0$ (using $n\cdot(m/n)
+ (m-n)\cdot0 = m$). This means $S=FF^*$ has $n$ eigenvalues all equal
to $m/n$ --- exactly the tight frame condition. The equiangular condition
is then encoded in the off-diagonal structure once tightness forces the
diagonal, so~\eqref{eq:etf-gram} captures both conditions simultaneously.
\end{proof}

\begin{corollary}[ETF necessary condition]
\label{cor:etf-necessary}
If an ETF with $m$ vectors in $\F^n$ exists, then the common coherence
value satisfies
\[
  c = \sqrt{\frac{m-n}{n(m-1)}},
\]
exactly the Welch bound.
\end{corollary}

\begin{proof}
If $\{f_i\}_{i=1}^m$ is an ETF, it achieves equality in the Welch bound
by Theorem~\ref{thm:welch-bound}. The Welch bound equals
$\sqrt{(m-n)/(n(m-1))}$, so the common inner product magnitude is exactly
that value.
\end{proof}

\section{The Grassmannian Packing Perspective}
\label{sec:grassmannian}

The Welch bound and ETFs arise naturally from a broader geometric
question: how should one arrange $m$ one-dimensional subspaces (lines
through the origin) in $\R^n$ (or $\C^n$) to make them as ``spread out''
as possible? A one-dimensional subspace in $\R^n$ is equivalent to a
pair of antipodal unit vectors $\pm f$ on the unit sphere. The set of
all such lines is called the \emph{real projective space} $\R\mathrm{P}^{n-1}$.
The question of packing $m$ lines in $\R^n$ to maximize their minimum
pairwise angle is called the \emph{Grassmannian packing problem}
(``Grassmannian'' because the set of all $k$-dimensional subspaces of
$\R^n$ is a manifold called the Grassmannian $\Gr(k,n)$, and the line
case $k=1$ is the simplest instance).

\begin{definition}[Grassmannian packing angle]
\label{def:packing-angle}
Given $m$ unit vectors $\{f_i\}$ in $\R^n$, the \emph{worst-case
packing angle} between the corresponding lines $\spn\{f_i\}$ is
\[
  \theta_{\min} :=
  \min_{i\ne j} \arccos\!\bigl(|\inner{f_i}{f_j}|\bigr)
  = \arccos\!\bigl(\mu(\{f_i\})\bigr).
\]
Since $\arccos$ is decreasing, \emph{maximizing $\theta_{\min}$} over
all unit-norm collections is equivalent to \emph{minimizing} the
coherence $\mu$.
\end{definition}

The Welch bound then says: no configuration of $m$ unit-norm vectors in
$\R^n$ can achieve $\theta_{\min}>\arccos\!\bigl(\sqrt{(m-n)/(n(m-1))}\bigr)$.
An ETF, if one exists for the given $(n,m)$ pair, is an optimal packing
--- a solution to the Grassmannian packing problem.

\begin{keyidea}
\textbf{ETFs are optimal Grassmannian packings}: they arrange $m$ lines
in $\R^n$ with maximum possible minimum angle between any pair. The
Welch bound is the ceiling on that minimum angle, expressed as a lower
bound on coherence. Tight frames minimize the frame potential (energy
spreading), while ETFs additionally maximize angular separation (packing
optimality). These are independent optimization criteria that happen to
be simultaneously achievable in certain cases, and those cases are
exactly the ETFs.
\end{keyidea}

\section{Harmonic Frames: ETFs from Roots of Unity}
\label{sec:harmonic-frames}

Chapter~4 (Lab Exercise~4.5) introduced the vectors
\[
  f_k = \frac{1}{\sqrt{n}}\bigl(1,\,\omega^k,\,\omega^{2k},\,\ldots,
  \omega^{(n-1)k}\bigr)\in\C^n, \qquad k=0,1,\ldots,n-1,
\]
where $\omega=e^{2\pi i/n}$, and observed that they form an orthonormal
basis for $\C^n$ (hence trivially a tight frame with $A=1$ and $\mu=0$).
\emph{Harmonic frames} generalise this construction to $m>n$ vectors by
selecting $m$ rows from the $n\times n$ DFT matrix --- choosing
$m$ indices $k$ from $\{0,1,\ldots,n-1\}$ rather than all $n$ of them.

\begin{definition}[Harmonic frame]
\label{def:harmonic-frame}
Let $\omega=e^{2\pi i/n}$ and let $\mathcal{S}\subset\Z/n\Z$ be a
subset of $m$ residues modulo $n$, with $m\le n$. The corresponding
\emph{harmonic frame} for $\C^n$ is
\[
  \bigl\{f_k\bigr\}_{k\in\mathcal S},\qquad
  f_k=\frac{1}{\sqrt{n}}
  \bigl(1,\,\omega^k,\,\omega^{2k},\,\ldots,\,\omega^{(n-1)k}\bigr)\in\C^n.
\]
\end{definition}

\begin{remark*}
This is most interesting when $m>n$ is larger than the dimension, which
requires going to a larger DFT. More generally, one takes an $n\times N$
DFT matrix (columns $\{f_k\}_{k=0}^{N-1}$ for some $N>n$) and selects
$m\le N$ columns; the discrete orthogonality relation for roots of unity
(Proof~2 of Theorem~4.3.1) guarantees tightness whenever the index set
$\mathcal S$ has a particular symmetry property (it suffices that
$\mathcal S$ is a ``difference set'' modulo $N$). We study only the
simplest case here.
\end{remark*}

\begin{proposition}[Harmonic frames from the full DFT are tight]
\label{prop:harmonic-tight}
The full DFT collection $\{f_k\}_{k=0}^{n-1}\subset\C^n$ defined above
is a Parseval frame for $\C^n$ (i.e.\ an orthonormal basis, $A=B=1$).
\end{proposition}

\begin{proof}
The $m\times m$ matrix $(F^*F)_{jk}=\inner{f_k}{f_j}=\frac1n\sum_{\ell=0}^{n-1}
\overline{\omega^{j\ell}}\omega^{k\ell}=\frac1n\sum_{\ell}\omega^{(k-j)\ell}$,
which equals $1$ when $j=k$ and $0$ when $j\ne k$ by the discrete
orthogonality relation for roots of unity (equation~(4.3) in Chapter~4,
introduced in Proof~2 of Theorem~4.3.1). So $F^*F=I_n$, confirming
the columns are orthonormal, hence a Parseval frame.
\end{proof}

The reason harmonic frames appear in a chapter on ETFs is that
\emph{selecting subsets} of DFT columns gives natural candidates for
equiangular collections, since the pairwise inner products
$\inner{f_j}{f_k}=(1/n)\sum_\ell\omega^{(k-j)\ell}$ depend only on
$k-j\pmod{n}$, not on the individual indices --- so if two pairs $(j_1,k_1)$
and $(j_2,k_2)$ have $k_1-j_1\equiv k_2-j_2\pmod n$, their inner products
are automatically equal. Whether \emph{all} off-diagonal inner products
end up equal (equiangularity) depends on the structure of $\mathcal S$,
and is related to the theory of \emph{difference sets} in combinatorics.

\begin{examplebox}[title={Example 9.4: A harmonic frame in $\C^3$}]
\label{ex:harmonic-c3}
Take $n=3$, $\omega=e^{2\pi i/3}$, and $\mathcal S=\{0,1,2\}$ (all three
columns). The harmonic frame is the orthonormal basis of $\C^3$:
\[
  f_0 = \tfrac{1}{\sqrt3}(1,1,1)^T,\quad
  f_1 = \tfrac{1}{\sqrt3}(1,\omega,\omega^2)^T,\quad
  f_2 = \tfrac{1}{\sqrt3}(1,\omega^2,\omega^4)^T.
\]
This is an ETF vacuously ($m=n=3$, so coherence is $0$ and the Welch
bound is $0$). The genuinely interesting case is taking $m<N$ columns
from a larger DFT; Lab Exercise~9.5 explores a specific example.
\end{examplebox}

\section{NumPy Implementation}
\label{sec:numpy-etf}

\begin{lstlisting}[caption={Coherence, Welch bound, and ETF verification in NumPy}]
import numpy as np

def gram_matrix(F):
    """Gram matrix G = F^* F (shape m x m)."""
    return F.T.conj() @ F

def coherence(F):
    """Maximum off-diagonal |inner product| for unit-norm columns of F."""
    G = gram_matrix(F)
    m = F.shape[1]
    off_diag = np.abs(G - np.diag(np.diag(G)))  # zero the diagonal
    return off_diag.max()

def welch_bound(n, m):
    """Lower bound on coherence for m unit-norm vectors in F^n."""
    return np.sqrt((m - n) / (n * (m - 1)))

def is_etf(F, tol=1e-6):
    """
    Check whether the columns of F form an equiangular tight frame.
    Returns (is_tight, is_equiangular, achieves_welch).
    """
    n, m = F.shape
    # Tightness
    S = F @ F.T.conj()
    eigs = np.linalg.eigvalsh(S.real)
    tight = bool(np.isclose(eigs.min(), eigs.max(), atol=tol))
    # Equiangularity: all off-diagonal |<fi,fj>| must be equal.
    # Compare max and min of off-diagonal magnitudes (not off[0],
    # which depends on ordering and may be zero for non-equiangular frames).
    G = gram_matrix(F)
    off = np.array([np.abs(G[i,j]) for i in range(m) for j in range(m) if i != j])
    equi = bool(np.isclose(off.max(), off.min(), atol=tol))
    # Welch bound
    mu = coherence(F)
    wb = welch_bound(n, m)
    achieves = bool(np.isclose(mu, wb, atol=tol))
    return tight, equi, achieves

# -------------------------------------------------------------------
# Mercedes-Benz frame (m=3, n=2): the canonical real ETF
# -------------------------------------------------------------------
f1 = np.array([1.0,  0.0])
f2 = np.array([-0.5, np.sqrt(3)/2])
f3 = np.array([-0.5, -np.sqrt(3)/2])
F_mb = np.column_stack([f1, f2, f3])

tight_mb, equi_mb, wb_mb = is_etf(F_mb)
print("=== Mercedes-Benz frame (m=3, n=2) ===")
print(f"Tight: {tight_mb}, Equiangular: {equi_mb}, Achieves Welch bound: {wb_mb}")
print(f"Coherence: {coherence(F_mb):.6f}  Welch bound: {welch_bound(2,3):.6f}")

# -------------------------------------------------------------------
# Square frame (m=4, n=2): tight but NOT equiangular
# -------------------------------------------------------------------
angles4 = np.array([0, np.pi/2, np.pi, 3*np.pi/2])
F_sq = np.vstack([np.cos(angles4), np.sin(angles4)])

tight_sq, equi_sq, wb_sq = is_etf(F_sq)
print("\n=== Square frame (m=4, n=2) ===")
print(f"Tight: {tight_sq}, Equiangular: {equi_sq}, Achieves Welch bound: {wb_sq}")
print(f"Coherence: {coherence(F_sq):.6f}  Welch bound: {welch_bound(2,4):.6f}")
G_sq = gram_matrix(F_sq)
print("Gram matrix off-diag magnitudes:",
      sorted(set(np.abs(G_sq - np.eye(4)).round(4).flatten()))[::-1][:4])

# -------------------------------------------------------------------
# Gram matrix check: G^2 = (m/n) G for Mercedes-Benz
# -------------------------------------------------------------------
G_mb = gram_matrix(F_mb)
lhs = G_mb @ G_mb
rhs = (3/2) * G_mb
print("\n=== Gram matrix check (Proposition 9.4.1) ===")
print(f"G^2 = (m/n)G for Mercedes-Benz: {np.allclose(lhs, rhs)}")

# -------------------------------------------------------------------
# Welch bound survey: all regular polygons m=3..8 in R^2
# -------------------------------------------------------------------
print("\n=== Regular m-gons vs Welch bound ===")
for m in range(3, 9):
    angles = 2*np.pi*np.arange(m)/m
    F = np.vstack([np.cos(angles), np.sin(angles)])
    mu = coherence(F)
    wb = welch_bound(2, m)
    t, e, w = is_etf(F)
    print(f"m={m}: mu={mu:.4f}, Welch={wb:.4f}, "
          f"tight={t}, equiangular={e}, ETF={t and e}")
\end{lstlisting}

\section{Searching for ETFs via Gradient Descent}
\label{sec:etf-search}

The frame potential (Chapter~8) finds tight frames but not necessarily
ETFs. To search for ETFs numerically, we can minimize a different
objective: the \emph{coherence itself}, or the \emph{coherence-squared
sum} $\sum_{i\ne j}|\inner{f_i}{f_j}|^2$. The latter is smoother (no
\texttt{max} operation) and directly related to the frame potential:
\[
  \sum_{i\ne j}|\inner{f_i}{f_j}|^2 = \FP - m \quad (\text{unit-norm case}).
\]
Since $\FP$ is minimized by tight frames, minimizing $\sum_{i\ne j}
|\inner{f_i}{f_j}|^2$ is the same optimization. This means that if an
ETF exists, gradient descent on the frame potential will find a
tight frame in which the off-diagonal inner products are likely to
equalize too, especially from a good random initialization.

\begin{examplebox}[title={Example 9.3: When gradient descent finds an ETF (and when it doesn't)}]
\label{ex:etf-gd}
For $(n,m)=(2,3)$: the only unit-norm tight frame in $\R^2$ with $3$
vectors is (up to rotation and sign-flips) the Mercedes-Benz frame, so
gradient descent on the frame potential automatically produces an ETF.

For $(n,m)=(2,4)$: there is no real ETF with $4$ vectors in $\R^2$ (as
the remark in Section~\ref{sec:coherence} explains), so gradient descent
finds tight frames that are not ETFs. The tight-frame minimum $m^2/n=8$
is achieved by many distinct non-equiangular configurations (the family
of tight non-equiangular frames is large), and gradient descent lands on
one of them.

For $(n,m)=(3,6)$: real ETFs with $6$ vectors in $\R^3$ are known to
exist (the regular icosahedron vertices, after normalization, give one
such configuration). Gradient descent can find them, but only if the
search is conducted over a refined objective that encourages
equiangularity beyond just tightness.
\end{examplebox}

\section{Chapter Summary}
\label{sec:summary-ch9}

\begin{itemize}
  \item The \textbf{coherence} $\mu=\max_{i\ne j}|\inner{f_i}{f_j}|$ of a
        unit-norm collection measures how ``non-orthogonal'' the most
        similar pair is. Small coherence means well-spread configurations.

  \item A collection is \textbf{equiangular} if all off-diagonal Gram
        entries have the same modulus. An \textbf{equiangular tight
        frame} (ETF) is equiangular and tight.

  \item \textbf{Theorem~\ref{thm:welch-bound} (Welch bound):} for any
        unit-norm collection of $m$ vectors in $\F^n$ with $m>n$,
        $\mu\ge\sqrt{(m-n)/(n(m-1))}$. Equality holds if and only if
        the collection is an ETF. The proof combines the frame-potential
        lower bound ($\FP\ge m^2/n$, from Chapter~8) with an upper bound
        on $\FP$ in terms of $\mu$.

  \item \textbf{ETFs are optimal Grassmannian packings}: they arrange $m$
        lines in $\R^n$ with the maximum possible minimum pairwise angle,
        and they minimize the frame potential simultaneously.

  \item \textbf{Proposition~\ref{prop:etf-gram}}: $\{f_i\}$ is an ETF
        iff its Gram matrix satisfies $G^2=(m/n)G$, a single matrix
        equation encoding both tightness and equiangularity.

  \item ETFs exist only for certain $(n,m)$ pairs. In $\R^2$, only $m=3$
        (the Mercedes-Benz frame) gives a real ETF. In higher dimensions
        and over $\C$, the existence problem is an active research area
        connected to combinatorics (e.g.\ symmetric designs, difference
        sets) and quantum information (optimal quantum measurements).
\end{itemize}

\section{Lab Exercises}
\label{sec:lab-ch9}

\begin{exercisebox}[title={Lab Exercise 9.1: Regular Polygons and Equiangularity}]
Implement \texttt{coherence(F)}, \texttt{welch\_bound(n, m)}, and
\texttt{is\_etf(F)} from Section~\ref{sec:numpy-etf}. Run
\texttt{is\_etf} on all regular $m$-gons in $\R^2$ for $m=3,\ldots,12$.
For each, also report the distinct off-diagonal Gram magnitudes. Confirm
that only $m=3$ is an ETF, and explain geometrically why even-$m$ cases
fail so dramatically (hint: $\mu=1$).
\end{exercisebox}

\begin{exercisebox}[title={Lab Exercise 9.2: Welch Bound Survey}]
For each $(n,m)$ in $\{(2,3),(2,5),(3,4),(3,6),(3,7),(4,6),(4,7)\}$:
\begin{enumerate}[label=(\alph*)]
  \item Compute the Welch bound $\sqrt{(m-n)/(n(m-1))}$.
  \item Use the gradient-descent algorithm of Chapter~8
        (\texttt{fit\_tight\_frame}) to find a tight unit-norm frame for
        each, and compute its coherence.
  \item Report whether the result achieves the Welch bound. For those
        that do not, compute the ratio $\mu/\mu_{\mathrm{Welch}}$ and
        call it the ``coherence excess.''
  \item Identify which $(n,m)$ pairs appear from your experiments to
        \emph{admit} an ETF (coherence excess $\approx 1$) and which do
        not (coherence excess $\gg 1$). (Note: this is a numerical
        experiment, not a proof --- you may miss some ETFs or
        misidentify some near-misses.)
\end{enumerate}
\end{exercisebox}

\begin{exercisebox}[title={Lab Exercise 9.3: The Gram Matrix Condition}]
For the Mercedes-Benz frame $F_\mathrm{MB}$, verify numerically that
$G^2=(3/2)G$ (Proposition~\ref{prop:etf-gram}). Then construct a
\emph{different} $3\times3$ Gram matrix $G'$ satisfying $G'^2=(3/2)G'$,
with $G'_{ii}=1$ for all $i$ and all off-diagonal entries having modulus
$\tfrac12$, but with the off-diagonal \emph{signs} permuted. Does $G'$
correspond to a valid frame? (Hint: a Gram matrix $G=F^*F$ must be
positive semidefinite; check whether $G'$ is PSD, and if so, extract the
corresponding frame vectors using a Cholesky decomposition or eigenvalue
decomposition.) This explores the phenomenon that multiple non-congruent
ETF configurations can share the same equiangularity value.
\end{exercisebox}

\begin{exercisebox}[title={Lab Exercise 9.4: Coherence-Minimizing Descent}]
The frame potential minimizes $\sum_{ij}|\inner{f_i}{f_j}|^2=\FP$.
Instead, minimize the \emph{maximum} off-diagonal inner product:
$\mu = \max_{i\ne j}|\inner{f_i}{f_j}|$, or a smooth approximation
such as $\mu_\alpha = \bigl(\sum_{i\ne j}|\inner{f_i}{f_j}|^{2\alpha}
\bigr)^{1/(2\alpha)}$ for large $\alpha$ (a soft-max approximation that
becomes exact as $\alpha\to\infty$). Implement gradient descent on
$\mu_\alpha$ with $\alpha=8$ for $(n,m)=(2,3)$ and $(n,m)=(3,4)$, and
compare the resulting coherence to the Welch bound. Does maximizing
angular separation produce results closer to the Welch bound than
minimizing the frame potential alone?
\end{exercisebox}

\begin{exercisebox}[title={Lab Exercise 9.5 (Optional, Challenge): ETFs in $\R^3$}]
The six vectors formed by taking the $\pm$ unit vectors along the three
coordinate axes of $\R^3$ --- i.e.\ $\{\pm e_1,\pm e_2,\pm e_3\}$ ---
form a tight frame for $\R^3$ with $m=6$, $A=2$. But is it equiangular?
(Hint: some pairs are orthogonal, others are antipodal.) If not, research
(via any source you can find) what the smallest real ETF in $\R^3$ with
$m>3$ is. Then verify the claimed ETF numerically using \texttt{is\_etf}.
This exercise previews the rich combinatorial theory of ETF existence,
which in $\R^3$ is already surprisingly non-trivial.
\end{exercisebox}

\begin{exercisebox}[title={Lab Exercise 9.6 (Optional): Harmonic Frames in $\C^n$}]
This exercise explores the harmonic frame construction of
Section~\ref{sec:harmonic-frames}.
\begin{enumerate}[label=(\alph*)]
  \item For $n=3$ and $\omega=e^{2\pi i/3}$, construct the three DFT
        vectors $f_0,f_1,f_2\in\C^3$ as in Definition~\ref{def:harmonic-frame}
        using complex NumPy arrays (\texttt{dtype=complex}). Verify
        numerically that $\inner{f_j}{f_k}=\delta_{jk}$ for all $j,k$
        (use \texttt{np.vdot} for complex inner products as discussed in
        Lab Exercise~1.3 of Chapter~1), confirming the orthonormal-basis
        claim of Proposition~\ref{prop:harmonic-tight}.
  \item Now take $N=6$, $\omega=e^{2\pi i/6}$, and consider all $6$ DFT
        vectors $\{f_k\}_{k=0}^5\subset\C^6$. Select the subset
        $\mathcal S=\{0,1,2\}$ (the first three), giving a collection of
        $3$ unit-norm vectors in $\C^6$. Is this collection tight? Is it
        equiangular? Compute its coherence and compare to the Welch bound
        for $n=6$, $m=3$.
  \item (Challenge) Investigate whether any size-3 subset of the $N=7$
        DFT columns gives an equiangular collection in $\C^7$. This is
        related to the existence of a $(7,3,1)$ difference set modulo $7$
        --- a set $\mathcal S$ such that every nonzero residue modulo $7$
        appears exactly once as a difference $s_1-s_2$ with $s_1\ne s_2$
        in $\mathcal S$. The set $\mathcal S=\{0,1,3\}$ is one such; test
        it numerically using \texttt{is\_etf} from Lab Exercise~9.1.
\end{enumerate}
\end{exercisebox}
\newpage
\section{Supplementary Exercises}
\label{sec:extra-ch9}

\subsection*{Theoretical Exercises}
\begin{theoexercises}
	\item Isolate ``Step~2'' of the Welch bound proof (Theorem~9.3.1) as a
	standalone lemma: prove that for any $m$ unit-norm vectors with
	coherence $\mu$, $\FP\le m+m(m-1)\mu^2$.
	
	\item Using $G^2=(m/n)G$ (Proposition~9.4.1), prove that the Gram
	matrix of an ETF has exactly two distinct eigenvalues: $0$ with
	multiplicity $m-n$, and $m/n$ with multiplicity $n$. Conclude
	$\rank G=n$.
	
	\item The \emph{chordal distance} between the lines spanned by unit
	vectors $f_i,f_j$ is $d(f_i,f_j):=\sqrt{1-\inner{f_i}{f_j}^2}$.
	Prove that minimizing the coherence $\mu$ over a configuration of
	$m$ unit vectors is equivalent to maximizing the \emph{minimum}
	pairwise chordal distance --- i.e.\ the Grassmannian packing
	problem (minimize worst-case inner product) and the
	``line-packing'' problem (maximize worst-case chordal distance)
	are the same optimization.
	
	\item State and prove the complex analogue of the Welch bound
	(Theorem~9.3.1) for unit vectors in $\C^n$, noting precisely where
	the proof requires the conjugate-linear inner product convention
	of Section~1.2.1.
	
	\item Numerically motivated conjecture, now to be proved: for the
	Mercedes-Benz ETF, show that $\sum_{i=1}^3\inner{x}{f_i}^4$ is the
	\emph{same} constant for every unit vector $x\in\R^2$ (not merely
	the second moment $\sum_i\inner{x}{f_i}^2$, which is constant for
	\emph{any} tight frame by definition). Then show this stronger
	``fourth-moment constancy'' \emph{fails} for the square
	configuration (tight, but not equiangular), so it is a genuinely
	finer invariant than tightness alone --- a first glimpse of what
	the frame theory literature calls a \emph{projective 2-design}.
\end{theoexercises}

\subsection*{Computational Exercises}
\begin{compexercises}
	\item Verify T1's bound numerically for $10$ random unit-norm
	configurations with $(n,m)=(3,8)$, confirming $\FP\le
	m+m(m-1)\mu^2$ in every case.
	
	\item Compute the Gram matrix eigenvalues of the Mercedes-Benz ETF and
	confirm they match T2's prediction ($0$ with multiplicity $m-n=1$,
	$m/n=3/2$ with multiplicity $n=2$).
	
	\item Implement the chordal distance function and verify T3's
	equivalence numerically: for several random configurations,
	confirm that the configuration minimizing coherence (via
	\texttt{fit\_min\_coherence} from Lab Exercise~9.4) also maximizes
	the minimum chordal distance.
	
	\item Reproduce the fourth-moment computation from T5: sample $50$
	unit vectors $x$ around the circle, compute
	$\sum_i\inner{x}{f_i}^4$ for both the Mercedes-Benz ETF and the
	square, and plot both as functions of the angle of $x$, visually
	confirming constancy for the ETF and variation for the square.
\end{compexercises}
